
\documentclass[10pt, a4paper, xcolor=dvipsnames]{beamer}

\usepackage{amsthm, graphicx}
%\usetheme{Berkeley}
\usetheme{Berlin}
\usefonttheme[onlylarge]{structuresmallcapsserif}
\usefonttheme[onlysmall]{structurebold}



\title[Neural Network Optimization]
{Neural Network Optimization Using Parallel Computing}
\subtitle{Project for Master's Degree in HPC}
\author{Saumya Bhatnagar}
\date{October 2, 2018}
\institute{Trinity College Dublin}

%\setlength{\unitlength}{1cm}
%\thicklines

\newcommand{\bs}{\textbackslash}


\usepackage{tikz}
\usetikzlibrary{tikzmark}

\newcounter{tmp}
\newcommand<>\Highlight[1]{%
	\stepcounter{tmp}%
	\only#2{\begin{tikzpicture}[remember picture,overlay]
		\fill[green!60!black,opacity=0.5] 
		([xshift=-.2em,yshift=2ex]pic cs:start-\thetmp)
		rectangle  
		([xshift=.2em,yshift=-1ex]pic cs:end-\thetmp);
		\end{tikzpicture}}%
	\tikzmark{start-\thetmp}#1\hfill\tikzmark{end-\thetmp}%
}

%\usepackage{biblatex}
%\setbeamertemplate{bibliography item}[text]
%\bibliographystyle{unsrtnat}

\usepackage[round,sort,comma,numbers]{natbib}
\usepackage[autostyle=true]{csquotes} % Required to generate language-dependent quotes in the bibliography



\usepackage{listings}
%\usepackage{minted}


\usepackage{smartdiagram}
\usetikzlibrary{shapes.geometric,calc}
\tikzset{
	planet/.append style={regular polygon, regular polygon sides=6},
	satellite/.append style={regular polygon, regular polygon sides=6},
	every picture/.append style={rotate=30},
	connection planet satellite/.style={
		bend right/.style=,
		every edge/.style={fill=\col},
		to path={
			\pgfextra
			\path[draw=none, fill=none] (\tikztostart) 
			-- coordinate[at start] (@start@) coordinate[at end] (@target@) (\tikztotarget);
			\endpgfextra
			\ifnum\xi<\maxsmitem % to disable the last arrow
			($(@start@)!.6cm!90:(@target@)$) -- ($(@target@)!.25cm!-90:(@start@)$)
			-- ($(@target@)!.25cm!90:(@start@)$) -- ($(@start@)!.6cm!-90:(@target@)$)
			-- cycle
			\fi}

			}
	}

\usepackage{smartdiagram}


\usetikzlibrary{positioning}


\begin{document}


\begin{frame}
\titlepage
\end{frame}

\section{Introduction \& Background}
\begin{frame}
%\frametitle{Essence of The Thesis}
\begin{itemize}
	\item<1-> Fourth Industrial Revolution \& Evolution of Machine Learning
	\item<2-> End of Moore's law and Quantum Mechanics\cite{MooreLaw}
	\item<3-> Evolution of HPC \cite{hennessy2011computer}
	\item<4-> Trends of HPC \cite{hennessy2011computer}
\end{itemize}
\hfill
\only<1>{
\begin{center}
	\resizebox{0.8\textwidth}{!}{
	\smartdiagramset{border color=none,
	back arrow disabled=true}
	\smartdiagram[flow diagram:horizontal]{
		\textbf{1950-Turing Test},
		1952- Self Game of Checkers,
		\textbf{1957-First Neural Network},
		1979-Mobile Robot,
		1985-Speech Recognition,
		\textbf{2006-DeepLearning}}}
\end{center}
}

\only<2> {\includegraphics[width=0.9\linewidth,height=5cm]{images/1_3.png}}
\only<3> {\includegraphics[width=0.8\linewidth,height=5cm]{images/1_4.jpg}}
\only<4> {\includegraphics[width=0.8\linewidth,height=5cm]{images/1_5.jpg}}
\end{frame}


\begin{frame}
\frametitle{Table of Contents}
\tableofcontents
\end{frame}

\section{Software Development}
\begin{frame}
%\frametitle{In This Thesis}
	\begin{minipage}{0.4\textwidth}
		\begin{enumerate}
			\item Compared accuracy of Neural Nets
			\item Single Population Fine Grained Algorithms
			\item Multiple-population Coarse Grained Algorithms\cite{cantu1998survey}
			\item Sequential Time Consumption
			\item Mpi4py, Keras, Tensorflow
			\item Analyzed TCHPC Clusters
			\item Memory Leak
		\end{enumerate}
	\end{minipage}
	\begin{minipage}{0.5\textwidth}
		\begin{centering}
			\resizebox{1\textwidth}{!}{
				\smartdiagram[connected constellation diagram]{
					Neural Network Optimization,
					Inspired by Genetic Algorithm,
					Quantum Mechanics,
					Parallel Computation,
					Python HPC,
					TCHPC Cluster,
					Cluster Nodes Comparison
				}
			}
		\end{centering}
	\end{minipage}
\end{frame}

\begin{frame}
%\frametitle{In This Thesis}
	\begin{itemize}
		\item To access best set of hyper-parameters 
		\item Convolutional Neural Network
		\item Image Processing
		\item Machine Learning Benchmark : CIFAR10
		\item Initiate with randomly selected combination of hyper-parameters
		\item Sample space kept same with every selection
		\item MPI parallelism
		\item Scheduling, Re-Scheduling by resuming the code based on Node availability using Slurm Workload Manager
		\item Keras has built-in support for multi-GPU data parallelism
		\item Boyle Cluster: ldd (GNU libc) 2.17
	\end{itemize} 
\end{frame}

\subsection{Optimization Method}
\begin{frame}
	\begin{minipage}{0.65\textwidth}
		\includegraphics[width=0.9\linewidth]{images/fig3b.png}
	\end{minipage}
	\begin{minipage}{0.25\textwidth}
		\begin{itemize}
			\item Inspired by GA \cite{hopfield1985neural}
			\item To crossbreed networks
			\item Relative accuracy of hyper-parameters set is analyzed
		\end{itemize}
	\end{minipage}
\end{frame}




\subsection{Parallel Computation}

\begin{frame}
%	\begin{minipage}{0.65\textwidth}
		\begin{center}
			\begin{tabular}{ c c c }
				Dataset & Data-points & Sequential Time (per generation) \\
				CIFAR10 & 10k & 15 minutes \\  
				80M Tiny Images & 80M & 20k minutes    
			\end{tabular}
		\end{center}
%	\end{minipage}
	Number of Processors\textless 32 =\textgreater Implemented Course Grained Parallelism \cite{li2007efficient}
	\newline
	\begin{minipage}{0.6\textwidth}

		Two Parallel Implementations of GA
		\begin{enumerate}
			\item<1-> \textbf{Single Population Fine Grained Algorithms} : using 4, 6, 8, 10, 12 processors
			\item<2-> \textbf{Multiple-population Coarse Grained Algorithms} : with inter-island(every 5 generations) and intra-island breeding
			\begin{itemize}
				\item<2-> population size of 4 in each island, using 12 processors
			\end{itemize}
		\end{enumerate}
	\end{minipage}
	\begin{minipage}{0.3\textwidth}
		\only<1> {\resizebox{\textwidth}{!}{\smartdiagram[circular diagram:clockwise]{1,2,3}}}
		\only<2> {\resizebox{\textwidth}{!}{
			\smartdiagram[circular diagram:clockwise]{	
				\resizebox{\textwidth}{!}{\smartdiagram[circular diagram:clockwise]{1,1,1}},
				\resizebox{\textwidth}{!}{\smartdiagram[circular diagram:clockwise]{2,2,2}},
				\resizebox{\textwidth}{!}{\smartdiagram[circular diagram:clockwise]{3, , }}
			}}}
	\end{minipage}
\end{frame}

\begin{frame}
\begin{columns}
	\begin{column}{0.68\textwidth}
		\includegraphics[width=0.85\linewidth]{images/fig3d.png}
	\end{column}
	\begin{column}{0.38\textwidth}
		\textbf{Single Population Fine Grained Algorithms} \\ 
		Send/Recv Model \\  
		One Network Per Node \\
		Inter-Network Crossover \newline \newline
		\textbf{Multiple-population Coarse Grained Algorithms} \\ 
		MPI Communicators \\  
		One Communicator Per Island \\
		Inter-Island \& Intra-Island Crossover
	\end{column}
\end{columns}
\end{frame}

\subsection{Implementation}
\begin{frame}[fragile]
	\textbf{MPI Initialization Testing}
	\begin{lstlisting}
MPI.Init()
print(MPI.Is_initialized())
print(MPI.Is_finalized())	
	\end{lstlisting}
	\textbf{Non-Blocking Exchange}
	\begin{lstlisting}
def nonBlockingExchange(self,data):
   reqSend1 = self.comm.isend(data,
  	  dest=((self.size+self.rank+1)\%self.size), 
 	  tag=self.rank)
   reqRecv2 = self.comm.irecv(source=
  	  ((self.size+self.rank-1)\%self.size), 
 	  tag=self.rank-1)
   dataPrev = reqRecv2.wait()
   reqSend1.wait()
   return dataPrev	
	\end{lstlisting}
\end{frame}

\begin{frame}[fragile]
	\textbf{MPI Gather}
	\begin{lstlisting}
recvdata = self.comm.Gather(data, root = 0)
	\end{lstlisting}
	\textbf{MPI Broadcast}
	\begin{lstlisting}
data = self.comm.bcast(data, root = 0)	
	\end{lstlisting}
	\textbf{Get Island}
	\begin{lstlisting}
self.subGroup = self.rank // groupSize
self.subComm = MPI.Comm.Split(
	self.comm, self.subGroup, self.rank)
	\end{lstlisting}
	\textbf{Get Island Details}
	\begin{lstlisting}
self.subSize, self.subRank = 
	self.subComm.Get_size(), self.subComm.Get_rank()
	\end{lstlisting}
\end{frame}


\subsection{Optimization Method - Quantum Inspired}
\begin{frame}[fragile]
Qubit in super-position state, increases sample space\cite{wang2005hybrid} \\
Hadamard gate
\begin{lstlisting}
r2=math.sqrt(2.0)           
h=np.array([[1/r2, 1/r2],[1/r2,-1/r2]])
\end{lstlisting}
Rotation Q-gate
\begin{lstlisting}
rot1=float(math.cos(theta));rot2=-float(math.sin(theta));
rot3=float(math.sin(theta));rot4=float(math.cos(theta));
\end{lstlisting}	
Quantum Population Vector
\begin{lstlisting}
# alpha squared
self.qpv[i,j,0]=np.around(2*pow(AlphaBeta[0],2),2) 
# beta squared
self.qpv[i,j,1]=1-self.qpv[i,j,0]
\end{lstlisting}
Quantum Measure
\end{frame}



\section{Analysis \& Conclusions}
\begin{frame}
	\begin{itemize}
		\item 	Search space consists of 5 hyper-parameters: activation functions, optimizers, hidden layers, nodes and dropout \cite{goodfellow2016deep}
		\item Six activation functions included are: sigmoid, elu, selu, relu, tanh, hard\_sigmoid
		 \item Six Optimizers include sgd, adagrad, adadelta, adam, adamax, nadam. 
		\item Hidden layers range from 1-15. 
		\item Nodes/neurons range from 4-128.
		\item Dropouts range from 0.1 to 0.5.
	\end{itemize}
\end{frame}


\subsection{Sequential Algorithm Analysis}
\begin{frame}
\frametitle{Samples tested}
\begin{columns}
	\begin{column}{0.48\textwidth}
		\includegraphics[width=0.9\linewidth]{../ResultAnalysis/optimizerVSActivation.jpg}
	\end{column}
	\begin{column}{0.58\textwidth}
		\begin{itemize}
			\item Activation functions: hard\_sigmoid and elu couldn't come in any set of combinations with the optimizer - adagrad
			\item Uneven Testing: some combinations have been tested more frequently \\
		\end{itemize}
	\end{column}
\end{columns}
\end{frame}

\begin{frame}
\frametitle{Samples tested}
\begin{columns}
	\begin{column}{0.54\textwidth}
		\includegraphics[width=\linewidth]{../ResultAnalysis/optimizerVSActivation.jpg}		\includegraphics[width=\linewidth]{../ResultAnalysis/neuronsVSlayers.jpg}
	\end{column}
	\begin{column}{0.54\textwidth}
		\begin{itemize}
			\item Activation functions: hard\_sigmoid and elu couldn't come in any set of combinations with the optimizer - adagrad
			\item Uneven Testing: some combinations have been tested more frequently \\
			\item Increasing the number of networks, increased the possible combinations
			\item Higher neurons and lower layers factored in high accuracies
		\end{itemize}
	\end{column}
\end{columns}
\end{frame}

\subsection{Sequential Algorithm vs Parallel Algorithm}
\begin{frame}
\frametitle{Fitness Comparison}
\begin{columns}
		\includegraphics[width=1.1\linewidth]{../ResultAnalysis/FitnessComparison_1.jpg}
\end{columns}
The convergence of parallel implementations is better than that of sequential implementation.
\end{frame}

\begin{frame}
\begin{columns}
	\includegraphics[width=1.2\linewidth]{../ResultAnalysis/FitnessComparison_2.jpg}
\end{columns}
\end{frame}

\begin{frame}
\frametitle{Time Comparison}
\begin{columns}
	\begin{column}{\textwidth}
		\includegraphics[width=\linewidth]{../ResultAnalysis/TimeComparison_1.jpg}
	\end{column}
\end{columns}
\begin{itemize}
	\item The time taken by sequential code increases with generations.
	\item Time taken by Sequential code is in hours while the same can be done in minutes using MPI parallelism
\end{itemize}
\end{frame}

\begin{frame}
\begin{columns}
	\begin{column}{1.1\textwidth}
		\includegraphics[width=\linewidth]{../ResultAnalysis/TimeComparison_2.jpg}
	\end{column}
\end{columns}
\end{frame}

\subsection{Parallel Algorithm - Island Model}
\begin{frame}
\frametitle{Multiple-population Coarse Grained Algorithms}
\begin{columns}
	\begin{column}{0.9\textwidth}
		\includegraphics[width=\linewidth]{../ResultAnalysis/TimeLapseIsland.pdf}
	\end{column}
	\begin{column}{0.2\textwidth}
		The Island model showed speed-up
	\end{column}
\end{columns}
\end{frame}

\subsection{Efficiency \& Speedups}
\begin{frame}
\frametitle{Super-linear behavior}
\begin{columns}
	\begin{column}{0.9\textwidth}
			\includegraphics[width=\linewidth]{../ResultAnalysis/Speedup_Efficiency.pdf}
	\end{column}
\end{columns}
\end{frame}

\begin{frame}
\frametitle{Super-linear behavior \& Memory leak}
\begin{columns}
	\begin{column}{0.48\textwidth}
		\includegraphics[width=\linewidth]{../ResultAnalysis/Speedup_Efficiency.pdf}
		Super-linear speedup
	\end{column}
	\begin{column}{0.58\textwidth}
		\includegraphics[width=\linewidth]{../ResultAnalysis/MemoryLeak.pdf}
	\end{column}
\end{columns}
\end{frame}


\begin{frame}
\frametitle{Memory Leak}
\begin{columns}
	\begin{column}{1.1\textwidth}
		\includegraphics[width=\textwidth]{../ResultAnalysis/MemoryLeak.jpg}
	\end{column}
\end{columns}
	Approx 100MB memory leak in 5 generations
\end{frame}




\subsection{Quantum Inspired Algorithm}
\frametitle{Parallel vs Quantum Inspired}
\begin{frame}
	\includegraphics[height=\textheight]{../ResultAnalysis/QEETO/OverallComparison.pdf}
\end{frame}

\begin{frame}
	\includegraphics[height=\textheight]{../ResultAnalysis/QEETO/Fitness.pdf}
\end{frame}


%\section{References}
\begin{frame}[allowframebreaks]
\frametitle{References}
%\scriptsize{\bibliographystyle{acm}}
\bibliographystyle{unsrtnat}
\bibliography{bibs/sample}
\end{frame}



\end{document}

